% GENERATED FILE. Edit canonical lessons, then run npm run generate:books.
% canonical-source-hash: fnv1a64:22d109235b63f4b8
% canonical-lessons: GU-C197-cokidar, GU-C197-tapasvu, GU-C197-olakhpatra, GU-C197-bandh, GU-C197-khulvu

\chapter{The watchman at the gate}
\label{ch:gu-the-watchman-at-the-gate}
\hlchaptermodality{197}

\begin{chapteropening}
\textbf{By the end of this chapter:} \emph{I can read opening times and show my card at the gate.}

Complete the last lesson of chapter 197: I can read opening times and show my card at the gate.
\end{chapteropening}

\section[\texorpdfstring{\gu{ચોકીદાર} (cokīdār)}{ચોકીદાર (cokīdār)}]{\gu{ચોકીદાર} (cokīdār) — watchman, guard}

\label{lesson:GU-C197-cokidar}



\begin{quote}
Before the new one: say the Gujarati for to go in, then the Gujarati for entry.
\end{quote}

\subsection*{You'll want to know: \gu{ચોકીદાર}}
\textbf{\gu{ચોકીદાર}} — \emph{cokīdār} — \textquotedblleft{}watchman, guard\textquotedblright{}.

\begin{quote}
\textbf{\gu{ચોકીદાર} \gu{રાત્રે} \gu{જાગે} \gu{છે}.} — \emph{cokīdār rātre jāge che} — \textquotedblleft{}The watchman stays awake at night.\textquotedblright{}
\end{quote}

\begin{cousinweb}
\textbf{\gu{ચોકીદાર}} (\emph{cokīdār}) is built from \emph{cokī}, \textquotedblleft{}guard post\textquotedblright{}, and the Persian suffix \emph{-dār}, \textquotedblleft{}holder, keeper\textquotedblright{}. CDIAL lists Gujarati \emph{cokī}, \textquotedblleft{}square low seat\textquotedblright{}, with \emph{cok}, \textquotedblleft{}square\textquotedblright{}, under Sanskrit \emph{catuṣka}, \textquotedblleft{}set of four\textquotedblright{} (CDIAL 4629).
\end{cousinweb}

\begin{grammarlens}[title={What you've built}]
One more word for everyday life at A2.
\end{grammarlens}

\subsection*{Your turn}
\begin{itemize}
\raggedright
  \item \emph{Say it:} \emph{cokīdār}
  \item \emph{Say it:} \emph{cokīdār}, once more
  \item \emph{Say it:} \emph{cokīdār rātre jāge che} — The watchman stays awake at night.
  \item \emph{Recall:} say the Gujarati for to go in, then the Gujarati for entry, then say \emph{cokīdār} again
\end{itemize}

\begin{tcolorbox}[breakable,colback=teal!4,colframe=teal!35!black,title={Before you move on}]
What does \gu{ચોકીદાર} mean? (\textquotedblleft{}Watchman, guard\textquotedblright{}.) Say it once more.
\end{tcolorbox}
\section[\texorpdfstring{\gu{તપાસવું} (tapāsvũ)}{તપાસવું (tapāsvũ)}]{\gu{તપાસવું} (tapāsvũ) — to check, to examine}

\label{lesson:GU-C197-tapasvu}



\begin{quote}
Before the new one: say the Gujarati for entry, then the Gujarati for watchman.
\end{quote}

\subsection*{You'll want to know: \gu{તપાસવું}}
\textbf{\gu{તપાસવું}} — \emph{tapāsvũ} — \textquotedblleft{}to check, to examine\textquotedblright{}.

\begin{quote}
\textbf{\gu{ચોકીદાર} \gu{દરેક} \gu{થેલી} \gu{તપાસે} \gu{છે}.} — \emph{cokīdār darek thelī tapāse che} — \textquotedblleft{}The watchman checks every bag.\textquotedblright{}
\end{quote}

\begin{cousinweb}
The origin of \textbf{\gu{તપાસવું}} (\emph{tapāsvũ}) is not settled. EtymDB links it to an Arabic verb for examining and lists Marathi \emph{tapāsṇe} beside it, but CDIAL has no entry for it, so the Arabic link is unconfirmed (EtymDB).
\end{cousinweb}

\begin{grammarlens}[title={What you've built}]
One more word for everyday life at A2.
\end{grammarlens}

\subsection*{Your turn}
\begin{itemize}
\raggedright
  \item \emph{Say it:} \emph{tapāsvũ}
  \item \emph{Say it:} \emph{tapāsvũ}, once more
  \item \emph{Say it:} \emph{cokīdār darek thelī tapāse che} — The watchman checks every bag.
  \item \emph{Recall:} say the Gujarati for entry, then the Gujarati for watchman, then say \emph{tapāsvũ} again
\end{itemize}

\begin{tcolorbox}[breakable,colback=teal!4,colframe=teal!35!black,title={Before you move on}]
What does \gu{તપાસવું} mean? (\textquotedblleft{}To check, to examine\textquotedblright{}.) Say it once more.
\end{tcolorbox}
\section[\texorpdfstring{\gu{ઓળખપત્ર} (oḷakhpatra)}{ઓળખપત્ર (oḷakhpatra)}]{\gu{ઓળખપત્ર} (oḷakhpatra) — identity card}

\label{lesson:GU-C197-olakhpatra}



\begin{quote}
Before the new one: say the Gujarati for watchman, then the Gujarati for to check.
\end{quote}

\subsection*{You'll want to know: \gu{ઓળખપત્ર}}
\textbf{\gu{ઓળખપત્ર}} — \emph{oḷakhpatra} — \textquotedblleft{}identity card\textquotedblright{}.

\begin{quote}
\textbf{\gu{ચોકીદારે} \gu{મારું} \gu{ઓળખપત્ર} \gu{જોયું}.} — \emph{cokīdāre mārũ oḷakhpatra joyũ} — \textquotedblleft{}The watchman saw my identity card.\textquotedblright{}
\end{quote}

A neuter noun.

\begin{cousinweb}
\textbf{\gu{ઓળખપત્ર}} (\emph{oḷakhpatra}) is built from \emph{oḷakh}, \textquotedblleft{}identity, recognition\textquotedblright{}, and \emph{patra}, \textquotedblleft{}leaf, letter\textquotedblright{}. CDIAL lists the verb \emph{oḷakhvũ}, \textquotedblleft{}to know, distinguish\textquotedblright{}, under Sanskrit \emph{upalakṣayati}, \textquotedblleft{}beholds\textquotedblright{} (CDIAL 2226), and \emph{patra} is Sanskrit \emph{páttra}, \textquotedblleft{}leaf\textquotedblright{} (CDIAL 7733).
\end{cousinweb}

\begin{grammarlens}[title={What you've built}]
One more word for everyday life at A2.
\end{grammarlens}

\subsection*{Your turn}
\begin{itemize}
\raggedright
  \item \emph{Say it:} \emph{oḷakhpatra}
  \item \emph{Say it:} \emph{oḷakhpatra}, once more
  \item \emph{Say it:} \emph{cokīdāre mārũ oḷakhpatra joyũ} — The watchman saw my identity card.
  \item \emph{Recall:} say the Gujarati for watchman, then the Gujarati for to check, then say \emph{oḷakhpatra} again
\end{itemize}

\begin{tcolorbox}[breakable,colback=teal!4,colframe=teal!35!black,title={Before you move on}]
What does \gu{ઓળખપત્ર} mean? (\textquotedblleft{}Identity card\textquotedblright{}.) Say it once more.
\end{tcolorbox}
\section[\texorpdfstring{\gu{બંધ} (bandh)}{બંધ (bandh)}]{\gu{બંધ} (bandh) — closed, shut}

\label{lesson:GU-C197-bandh}



\begin{quote}
Before the new one: say the Gujarati for to check, then the Gujarati for identity card.
\end{quote}

\subsection*{You'll want to know: \gu{બંધ}}
\textbf{\gu{બંધ}} — \emph{bandh} — \textquotedblleft{}closed, shut\textquotedblright{}.

\begin{quote}
\textbf{\gu{દુકાન} \gu{રવિવારે} \gu{બંધ} \gu{છે}.} — \emph{dukān ravivāre bandh che} — \textquotedblleft{}The shop is closed on Sundays.\textquotedblright{}
\end{quote}

\begin{cousinweb}
\textbf{\gu{બંધ}} (\emph{bandh}) is Sanskrit \emph{bandhá}, \textquotedblleft{}bond, damming, custody\textquotedblright{}, taken over whole, a tatsama; CDIAL lists its inherited cousin \emph{bā̃dh}, \textquotedblleft{}dam\textquotedblright{} (CDIAL 9136). CDIAL does not give the sense \textquotedblleft{}shut\textquotedblright{}, which also matches Persian \emph{band}, and the two words may have run together.
\end{cousinweb}

\begin{grammarlens}[title={What you've built}]
One more word for everyday life at A2.
\end{grammarlens}

\subsection*{Your turn}
\begin{itemize}
\raggedright
  \item \emph{Say it:} \emph{bandh}
  \item \emph{Say it:} \emph{bandh}, once more
  \item \emph{Say it:} \emph{dukān ravivāre bandh che} — The shop is closed on Sundays.
  \item \emph{Recall:} say the Gujarati for to check, then the Gujarati for identity card, then say \emph{bandh} again
\end{itemize}

\begin{tcolorbox}[breakable,colback=teal!4,colframe=teal!35!black,title={Before you move on}]
What does \gu{બંધ} mean? (\textquotedblleft{}Closed, shut\textquotedblright{}.) Say it once more.
\end{tcolorbox}
\section[\texorpdfstring{\gu{ખૂલવું} (khūlvũ)}{ખૂલવું (khūlvũ)}]{\gu{ખૂલવું} (khūlvũ) — to open up, to come open}

\label{lesson:GU-C197-khulvu}



\begin{quote}
Before the new one: say the Gujarati for identity card, then the Gujarati for closed.
\end{quote}

\subsection*{You'll want to know: \gu{ખૂલવું}}
\textbf{\gu{ખૂલવું}} — \emph{khūlvũ} — \textquotedblleft{}to open up, to come open\textquotedblright{}.

\begin{quote}
\textbf{\gu{બેંક} \gu{સવારે} \gu{ખૂલે} \gu{છે}.} — \emph{bẽk savāre khūle che} — \textquotedblleft{}The bank opens in the morning.\textquotedblright{}
\end{quote}

\textbf{\gu{ખોલવું}} opens a thing; \textbf{\gu{ખૂલવું}} is the thing opening.

\begin{cousinweb}
\textbf{\gu{ખૂલવું}} (\emph{khūlvũ}) is a tadbhava. CDIAL lists it, glossed \textquotedblleft{}to be open\textquotedblright{}, under \emph{*khull}, \textquotedblleft{}to be open\textquotedblright{}, a reconstructed form; the taught \emph{kholvũ}, \textquotedblleft{}to open\textquotedblright{}, is its partner from \emph{*khōll} (CDIAL 3945).
\end{cousinweb}

\begin{grammarlens}[title={What you've built}]
One more word for everyday life at A2.
\end{grammarlens}

\subsection*{Your turn}
\begin{itemize}
\raggedright
  \item \emph{Say it:} \emph{khūlvũ}
  \item \emph{Say it:} \emph{khūlvũ}, once more
  \item \emph{Say it:} \emph{bẽk savāre khūle che} — The bank opens in the morning.
  \item \emph{Recall:} say the Gujarati for identity card, then the Gujarati for closed, then say \emph{khūlvũ} again
\end{itemize}

\begin{tcolorbox}[breakable,colback=teal!4,colframe=teal!35!black,title={Before you move on}]
What does \gu{ખૂલવું} mean? (\textquotedblleft{}To open up, to come open\textquotedblright{}.) Say it once more.
\end{tcolorbox}
